%% packages

% typesetting
\usepackage[utf8]{inputenc}
\usepackage[T1]{fontenc}
\usepackage{microtype}
\usepackage{tikz-cd}
\usepackage{booktabs}
\usepackage{enumitem}
%\usepackage{parskip}

\usepackage[margin=1.4in]{geometry}

% mathematics
\usepackage{mathtools}
\usepackage{amsthm}
\usepackage{amssymb}
\usepackage{mathrsfs}

% references
\usepackage[sorting=nyt,maxnames=4]{biblatex}
\usepackage{hyperref}
\hypersetup{
    colorlinks,
    linkcolor={red!50!black},
    citecolor={blue!70!black},
    urlcolor={blue!70!black}
}
\usepackage{footnotebackref}
\usepackage[capitalise]{cleveref}

%% commands

% theorems
\newtheorem{theorem}[equation]{Theorem}
\newtheorem{proposition}[equation]{Proposition}
\newtheorem{lemma}[equation]{Lemma}
\newtheorem{corollary}[equation]{Corollary}

\theoremstyle{definition}
\newtheorem{definition}[equation]{Definition}

\theoremstyle{remark}
\newtheorem{remark}[equation]{Remark}
\newtheorem{example}[equation]{Example}

% references
\numberwithin{equation}{section}
\numberwithin{figure}{section}
\crefname{section}{\S\kern -3pt}{\S\S}
\crefname{equation}{}{}

% typesetting
\renewcommand*{\bibfont}{\normalfont\footnotesize}
\let\oldtocsection=\tocsection
\let\oldtocsubsection=\tocsubsection
\let\oldtocsubsubsection=\tocsubsubsection
\renewcommand{\tocsection}[3]{\hspace{0em}\oldtocsection{#1}{#2}{#3}}
\renewcommand{\tocsubsection}[3]{\hspace{2.5em}{\oldtocsubsection{#1}{#2}{#3}}}

% operators
\DeclareMathOperator{\Spec}{Spec}
\DeclareMathOperator{\Tot}{Tot}
\DeclareMathOperator{\cone}{cone}
\DeclareMathOperator{\Hom}{Hom}
\DeclareMathOperator{\Ext}{Ext}
\DeclareMathOperator{\shHom}{\mathcal{H}\kern -1pt \textit{om}}
\DeclareMathOperator{\shExt}{\mathcal{E}\kern -1pt \textit{xt}}
\DeclareMathOperator{\End}{End}
\DeclareMathOperator{\Crit}{Crit}
\DeclareMathOperator{\RHom}{\dR Hom}
\DeclareMathOperator{\RG}{\dR\Gamma}
\DeclareMathOperator{\Perf}{\mathfrak{Perf}}
\DeclareMathOperator{\Bl}{Bl}

% macros
\newcommand{\lperp}[1]{\prescript{\perp}{}{#1}} % left orthogonal
\newcommand{\todo}[1]{{\color{red}#1}}

% names
\newcommand{\id}{\mathrm{id}}
\newcommand{\op}{\mathrm{op}}
\newcommand{\pt}{\mathrm{pt}}
\newcommand{\dR}{\mathrm{R}} % right derived prefix R

% letters
\newcommand{\Z}{\mathbb{Z}}
\newcommand{\C}{\mathbb{C}}
\newcommand{\A}{\mathbb{A}}
\newcommand{\scrC}{\mathscr{C}}
\renewcommand{\P}{\mathbb{P}}
\renewcommand{\O}{\mathcal{O}}
